\chapter{Architecture Decision Records}
\label{ch:adr}

This chapter presents the Architecture Decision Records (ADRs) for Part~2. ADRs 1 to 4 are carried over from Part~1 with implementation notes. ADRs 5 and 6 are new decisions made during implementation.

\section*{ADR-1: Transactional Outbox (carried over, status: Implemented)}
\addcontentsline{toc}{section}{ADR-1: Transactional Outbox}

\textbf{Status:} Accepted. Implemented in \texttt{OmnichannelOutbox} plugin.

\subsection*{Decision}
Write cross-context events to an \texttt{OutboxMessage} table inside the same DB transaction as the order, then dispatch to RabbitMQ asynchronously via a polling \texttt{IScheduleTask}.

\subsection*{Implementation notes}
The \texttt{TransactionScopeAsyncFlowOption.Enabled} constraint discovered in the feasibility spike was confirmed in production code. Isolation level was tuned to \texttt{ReadCommitted} as anticipated. No further divergence from the Part~1 decision.

\subsection*{Alternatives considered (Part~1)}
\begin{itemize}
    \item Direct RabbitMQ publish from \texttt{IConsumer<OrderPlacedEvent>}: REJECTED (dual-write risk).
    \item XA / two-phase commit: REJECTED (RabbitMQ not an XA resource; operational complexity).
    \item CDC / Debezium log mining: REJECTED (engine-specific; out of scope).
    \item Synchronous publish inside the order transaction: REJECTED (broker in critical path).
\end{itemize}

\section*{ADR-2: Optimistic Stock Reservation (carried over, status: Implemented)}
\addcontentsline{toc}{section}{ADR-2: Optimistic Stock Reservation}

\textbf{Status:} Accepted. Implemented via \texttt{OrderPlacedEvent} outbox path and \texttt{StockReservedHandler} / \texttt{ReservationRejectedHandler}.

\subsection*{Decision}
Accept orders optimistically without a synchronous stock check. The \texttt{OrderPlaced} event instructs the Inventory service to attempt a reservation. If the reservation fails, the compensation flow transitions the order to \texttt{Compensated}.

\subsection*{Implementation notes}
The \texttt{Compensated} state is now concrete in \texttt{OrderStatus}. The compensation flow is implemented in \texttt{CompensationService}. No divergence from the Part~1 decision.

\section*{ADR-3: WireMock Carrier Simulation (carried over, status: Implemented)}
\addcontentsline{toc}{section}{ADR-3: WireMock Carrier Simulation}

\textbf{Status:} Accepted. Implemented in the Shipping Service.

\subsection*{Decision}
Use WireMock to simulate the carrier API, avoiding a dependency on a live external service during development and the demo.

\section*{ADR-4: Selective Context Extraction (carried over, status: Implemented)}
\addcontentsline{toc}{section}{ADR-4: Selective Context Extraction}

\textbf{Status:} Accepted. Inventory and Shipping services extracted; Catalog, Order Management, and Identity remain in the monolith.

\subsection*{Decision}
Extract only the contexts whose failure modes are operationally distinct (Inventory, Shipping). Keep the remaining contexts inside the nopCommerce monolith to avoid unnecessary operational complexity for a demo-scale system.

\section*{ADR-5: Two Plugins Instead of One (new)}
\addcontentsline{toc}{section}{ADR-5: Two Plugins Instead of One}

\textbf{Status:} Accepted. Decided during Part~2 implementation.

\subsection*{Context}
Part~1 described a single \texttt{OmnichannelOutbox} plugin. During implementation the outbox write path and the inbound subscriber path grew in different directions, with distinct failure domains and lifecycle concerns.

\subsection*{Decision}
Split into two plugins: \texttt{OmnichannelOutbox} (write path) and \texttt{OmnichannelConsumers} (read/inbound path).

\subsection*{Rationale}
\begin{itemize}
    \item Each plugin has a single deployment and failure mode.
    \item The two paths can be enabled/disabled independently without affecting the other.
    \item Mirrors the ADD iteration structure: Iteration~1 = outbox write, Iteration~2 = consumers.
\end{itemize}

\subsection*{Alternatives considered}
\begin{itemize}
    \item \textbf{Single plugin:} REJECTED: one plugin holding both the dispatcher and the subscriber creates a large, hard-to-test unit with mixed failure domains.
    \item \textbf{Separate process for consumers:} REJECTED: out of scope; the monolith process is already hosting the outbox dispatcher as an \texttt{IScheduleTask}.
\end{itemize}

\section*{ADR-6: Inbox Table Deduplication (new)}
\addcontentsline{toc}{section}{ADR-6: Inbox Table Deduplication}

\textbf{Status:} Accepted. Implemented in \texttt{InboxDeduplicationService} and \texttt{ProcessedInboxMessage} entity.

\subsection*{Context}
RabbitMQ delivers at-least-once. Part~1 stated that consumers must be idempotent but left the mechanism open.

\subsection*{Decision}
Persist a \texttt{ProcessedInboxMessage} record (keyed on \texttt{EventId}) for every handled message. Check for this record before processing. Duplicates are discarded without executing handler logic.

\subsection*{Rationale}
\begin{itemize}
    \item Handler logic (state transitions, payment voids, notifications) is not naturally idempotent; centralising the check at the infrastructure layer removes the burden from individual handlers.
    \item The \texttt{ProcessedInboxMessage} table doubles as an audit log of processed messages.
    \item The same \texttt{InboxRetentionTask} pattern used by the outbox cleans up old records.
\end{itemize}

\subsection*{Alternatives considered}
\begin{itemize}
    \item \textbf{Handler-level idempotency:} REJECTED: each handler would need to query its own state to determine whether the action had already been taken (e.g., ``is this order already Confirmed?''). This is error-prone and forces domain logic to carry infrastructure concerns.
    \item \textbf{Exactly-once delivery via RabbitMQ quorum queues:} REJECTED: quorum queues reduce but do not eliminate duplicates; does not remove the need for consumer-side idempotency.
\end{itemize}
